\section{Composition}\label{sec:composition}

\subsection{Model}\label{sec:model}

A fixed set~$V$ of $n\geq1$ validators is given.  Every validator has
weight~$1$: for $S\subseteq V$, $w(S)=|S|$, and the quorum threshold is
$q=\lceil 2n/3\rceil$.  $\Delta$ is the delivery parameter of blocks and
AC messages, and $\Delta_F\geq\Delta$ is the delivery parameter of
attestations.  A bound is strict: the receiver processes the object at
a time strictly less than the deadline of the bound.  When the deadline is
a public time, the receiver thus processes the object before the tick at
that time.  Time is continuous, and the clock~$t$ of a node is the
current time.  The public times are the multiples of~$\Delta$.  A node
runs one tick at each public time.  A tick starts from the state as it
was before its public time, and each step of the tick reads the state
after the earlier steps of that tick.

A node is awake or asleep.  An asleep node receives nothing and runs no
tick, so it has no duty and no output.  An object that is sent while its
receiver sleeps, or that is not yet delivered when its receiver goes to
sleep, is delivered after the next wake: a delivery bound counts from the
later of the send time and that wake.  The node stamps an object with its
time of receipt.  A node that wakes restarts (Section~\ref{sec:store})
before it runs a tick or processes an object at that time; a crash is a
sleep.

The delivery bounds hold for every awake honest node.  An honest node
forwards an object when it accepts it: a block when the FG accepts it, an
AC message when the AC admits it, and an attestation when it passes the
checks.  This includes the objects that a block carries: a vote that the AC
admits from the payload of an accepted block, and an attestation of an
accepted block that enters the pool, are forwarded as wire objects of their
own, whether or not another node accepts the block.  So a delivery bound also counts from the time at which an honest
node first accepts the object; a receipt that is refused is not forwarded.  A node that runs
while its delivery is under the control of the adversary is faulty.

Slot~$s$ starts at the public time~$t_s$, with $t_0=0$ and
$t_s<t_{s+1}$, and has an assigned proposer.  The slot of a time~$t$, written $\slot(t)$, is
the greatest~$s$ with $t_s\leq t$.  Round~$r$ has the slots
$rR,\ldots,rR+R-1$, and $\round(s)=\lfloor s/R\rfloor$.  Its opening time
is $T_r=t_{rR}$.  Rounds are the heartbeat of the network for the
aggregation of attestations; no rule of the shared node reads them, and an
AC can use them.  The AC fixes the slot times and a lower bound on~$R$.

The local validator index is~$\valindex$.

\subsection{Blocks}\label{sec:blocks}

Every block~$B$ has a root~$B.\rootf$, slot~$B.\slot$, parent~$B.\parent$,
and proposer index~$B.\proposer$.  It has an AC payload~$B.\acp$ and a list
$B.\mathsf{attestations}$ of attestations (the FG payload).  The
order of each list and its repeated entries are part of the block.  The
block names its parent by root; the FG state must hold that parent in full,
with the same content, not only a block with the same root.  $B.\rootf$
is the fixed digest of
\[
 (B.\parent.\rootf,\ B.\slot,\ B.\mathsf{attestations},\ B.\acp).
\]
The AC fixes the encoding of its payload in this input.

The proposer signs a block with the slot key of $B.\proposer$.  A
validator signs an attestation with its fixed identity $\valindex\in V$.
A handler receives an object only when its signature is valid for its
signer and these checks pass: the signer of an attestation is in~$V$ and
its pairs have the forms of Section~\ref{sec:fg}; the slot key of
$B.\proposer$ signs a block~$B$; every attestation
that $B$ carries passes the same checks as on the wire;
$\Call{ac\_valid\_payload}{B}$ holds; an AC message passes the checks of
the AC; and $B.\rootf$ is the digest of the content of~$B$.  The check
drops an object that fails it, with no effect.

Processed blocks form a parent-closed tree rooted at~$\genesis$.  Write
$B\preceq C$ when $B=C$ or $B$ is an ancestor of~$C$, and write $B\prec C$
for strict ancestry.  Two blocks are compatible when one is an ancestor of
the other, and they conflict otherwise.  The depth of a block is the
number of parent edges from genesis.  ``Deepest'' means maximum depth;
ties use a fixed root order.

\subsection{State}\label{sec:store}

The node keeps the clock~$t$, the FG state~$\Sigma$ and the AC
state~$\AC$.  The FG state is
\[
 \Sigma=(\T,\sigma[\cdot],F,J,h_j,\hmax,\attestations).
\]
$\Sigma.\T$ is the tree of accepted blocks.  A block enters it only when
it extends~$\Sigma.F$, and no rule removes a block.
$\Sigma.\sigma[B]$ is the immutable chain state of an accepted
block.  $\Sigma.F$, $\Sigma.J$, $\Sigma.h_j$ and $\Sigma.\hmax$ are the
finality caches.  Section~\ref{sec:fg} defines the chain state and the
caches.  $\Sigma.\attestations$ holds
the processed attestations in full, one copy of each, in the order of
first acceptance, with no age limit and no limit per validator.  No rule
reads the time at which an attestation arrived; the pool serves inclusion
only.  Attestations reach the pool on the wire, through the node's own
$\Call{attest}{}$, and from the blocks that the FG accepts, after the
block's finality update.

Only the FG functions read and write~$\Sigma$, and only the AC functions
read and write~$\AC$.  The AC has no block store: the FG decides which
blocks are accepted, and the AC sees them through the FG frame.

The validator also keeps its anti-slashing record~$\Lambda$
(Section~\ref{sec:fg}).  The record belongs to the validator.  It is
durable: it persists through a restart.  A restart keeps~$\Lambda$,
$\Sigma$ and the clock, and it replaces~$\AC$ by
$\Call{ac\_restart}{\AC,\Fr,t}$, with $\Fr=\Call{fg\_to\_ac}{\Sigma}$
and $t$ the time of the restart.  The AC decides which part of its state
persists.

At initialization, $t=0$,
\[
 \Sigma.\T=\{\genesis\},\quad
 \Sigma.F=\Sigma.J=\genesis,\quad \Sigma.h_j=0,\quad \Sigma.\hmax=1,
\]
$\Sigma.\sigma[\genesis]$ is the initial chain state, the pool is empty,
and $\AC$ is the initial state of the AC.

\subsection{The FG frame}\label{sec:frame}

The FG frame is what the FG gives to the AC.  It is derived from the FG
state when used, and every call of one tick step reads the FG frame of the
same state:
\[
 \Fr=\Call{fg\_to\_ac}{\Sigma}=(\tree,\fcroot,\allowed).
\]
\begin{itemize}
\item $\Fr.\tree$ is the set of the accepted blocks that
  extend~$\Sigma.F$.  Write $\Fr.F$ for its root, the block~$\Sigma.F$.
\item $\Fr.\fcroot$ is a block of $\Fr.\tree$: the root of the fork
  choice.
\item $\Fr.\allowed$ is a set of blocks of $\Fr.\tree$ that
  extend~$\Fr.\fcroot$: the blocks on which a head walk may stop.
\end{itemize}
For a set $\mathit{tree}\subseteq\Fr.\tree$, the \emph{viable tree}
$\filtered(\Fr,\mathit{tree})$ is the set of the blocks of~$\mathit{tree}$
that extend~$\Fr.\fcroot$ and have a descendant in
$\mathit{tree}\cap\Fr.\allowed$.  It follows from the FG frame; it is not a
part of it.
$\Fr.F$ moves only to a descendant, also through a restart, because a
restart keeps~$\Sigma$.  So a block that leaves $\Fr.\tree$ does not
return.  Figure~\ref{alg:fg-store} computes the FG frame.

\subsection{The available chain}\label{sec:ac-interface}

An AC has a state~$\AC$, a set of wire messages, a block payload, and
these functions.  $t$ is the clock of the node.
\[
\begin{array}{l@{\qquad}l}
 \Call{ac\_valid\_payload}{B} & \text{stateless check of $B.\acp$},\\
 \Call{ac\_on\_message}{\AC,\Fr,t,m} & \text{an AC wire message arrives},\\
 \Call{ac\_on\_block}{\AC,\Fr,t,B} & \text{the FG accepted block $B$},\\
 \Call{ac\_on\_tick}{\AC,\Fr,t} & \text{the schedule of the AC; it can
   send AC messages},\\
 \Call{ac\_proposal}{\AC,\Fr,t} & \text{returns
   $(H,\mathit{ac\_payload})$: a parent and an AC payload, or $\bot$},\\
 \Call{ac\_to\_fg}{\AC,\Fr,t} & \text{returns
   $(\mathit{fg\_confirmed},\mathit{head})$; $\mathit{fg\_confirmed}=\bot$
   when no block is new},\\
 \Call{ac\_confirmed}{\AC,\Fr,t} & \text{returns the confirmed tip},\\
 \Call{ac\_head}{\AC,\Fr,t} & \text{returns the head},\\
 \Call{ac\_restart}{\AC,\Fr,t} & \text{returns the AC state after a restart}.
\end{array}
\]
The two receipt handlers $\Call{ac\_on\_message}{}$ and
$\Call{ac\_on\_block}{}$ can also send AC messages, as
$\Call{ac\_on\_tick}{}$ can.  These functions do not take the FG state.  The AC holds a block in its state only as
a name: the content of a block and its ancestry come from $\Fr.\tree$, or
from the block that a call supplies.  The AC must satisfy these
requirements.
\begin{enumerate}
\item Every head of a duty (the parent of a proposal, the head of a vote,
  the head that $\Call{ac\_to\_fg}{}$ returns) extends~$\Fr.\fcroot$.
  Such a head whose walk starts in $\filtered(\Fr,\mathit{tree})$ is in
  $\Fr.\allowed$.
\item A block $\mathit{fg\_confirmed}$ other than~$\bot$ is in
  $\filtered(\Fr,\Fr.\tree)$, and the returned head extends it.
\item The output of $\Call{ac\_head}{}$ extends~$\Fr.\fcroot$, and the
  confirmed tip extends~$\Fr.F$ and is an ancestor of that output.
\item The parent of a proposal is a head of the AC at that time.
\end{enumerate}

The protocol exposes three tips: the finalized tip~$\Sigma.F$, the
confirmed tip $\Call{ac\_confirmed}{\AC,\Fr,t}$ and the head
$\Call{ac\_head}{\AC,\Fr,t}$, with $\Fr=\Call{fg\_to\_ac}{\Sigma}$.  They
lie on one chain: the finalized tip is an ancestor of the confirmed tip,
and the confirmed tip is an ancestor of the head.

\subsection{Node}\label{sec:node}

The node runs the tick of a public time before it processes any object
received at that time.  It stamps an object with the time at which it
processes it, so an object received at a public time~$\tau$ counts as
received after the tick at~$\tau$.  A tick runs the schedule of the AC,
then the proposal duty at~$t_s$,
then the FG update.  A node processes its own block, its own AC
message and its own attestation at once.

The node drops a block of a future slot.  The FG then accepts the block
only when its full parent is held, no held block has its root, it
extends~$\Sigma.F$, and the assertions of its state transition hold.  A
failed test or assertion refuses the block before any write, and the AC
does not see a refused block.
No handler then runs for an object that it carries, and a later delivery
can call the handler again.  The proposer takes the parent and the AC
payload from the AC; when the AC returns~$\bot$, it makes no proposal.  It includes every attestation of its pool that is
valid against the state of the parent and is not already on the parent's
chain, in the order of its pool.

\emph{The FG update.}  After each tick the node asks the AC for an FG
update.  The AC returns a pair $(\mathit{fg\_confirmed},\mathit{head})$, with
$\mathit{fg\_confirmed}\neq\bot$ only when it has a new block to vouch
for.
$\mathit{fg\_confirmed}$ is the block up to which the AC vouches for the
chain of the head, for the FG.  The node gives the FG the block
$\mathit{fg\_confirmed}$ of each such pair, the \emph{source}, and the
client record decides whether an attestation follows: a target only as the
first vote at a height, and no entry twice (Section~\ref{sec:fg}).  The
height pair comes from the chain state of the source, and the finality
pair from the entry state of that chain state
(Section~\ref{sec:fg-client}).  The FG does not read the head.

The node functions take the input from outside the node: the time of a
tick, a received object, or the output of the AC for a duty.  They read
the clock~$t$ and pass $\Sigma$, $\AC$ and~$\Lambda$ to the two parts.

\captionof{figure}{Node: handlers and duties}\label{alg:node}
\nopagebreak
{\footnotesize
\begin{algorithmic}[1]
\Function{on\_tick}{$t$}
  \Comment{at each public time $t$}
  \State $s\gets\slot(t)$
  \State $\Fr\gets\Call{fg\_to\_ac}{\Sigma}$; \Call{ac\_on\_tick}{$\AC,\Fr,t$}
  \If{$s>0$ and $t=t_s$ and $\valindex$ is the slot-$s$ proposer}
    \State $P\gets\Call{ac\_proposal}{\AC,\Fr,t}$
    \Comment{$\bot$: the AC makes no proposal}
    \If{$P\neq\bot$} \State \Call{propose\_block}{$P$} \EndIf
  \EndIf
  \State $\Fr\gets\Call{fg\_to\_ac}{\Sigma}$;\quad
    $(\mathit{fg\_confirmed},\cdot)\gets\Call{ac\_to\_fg}{\AC,\Fr,t}$
  \Comment{$\bot$: no new block for the FG}
  \If{$\mathit{fg\_confirmed}\neq\bot$}
    \State \Call{attest}{$\mathit{fg\_confirmed}$}
    \Comment{the source; the FG does not read the head}
  \EndIf
\EndFunction
\vspace{1pt}
\Function{on\_block}{$B$}
  \If{$B.\slot>\slot(t)$} \State \Return \EndIf
  \If{not $\Call{fg\_on\_block}{\Sigma,B}$} \State \Return
    \Comment{the FG accepts or refuses the block}
  \EndIf
  \State $\Fr\gets\Call{fg\_to\_ac}{\Sigma}$; \Call{ac\_on\_block}{$\AC,\Fr,t,B$}
    \Comment{the frame after the FG accepted $B$}
\EndFunction
\vspace{1pt}
\Function{on\_ac\_message}{$m$}
  \State $\Fr\gets\Call{fg\_to\_ac}{\Sigma}$; \Call{ac\_on\_message}{$\AC,\Fr,t,m$}
\EndFunction
\vspace{1pt}
\Function{on\_attestation}{$a$}
  \State \Call{fg\_on\_attestation}{$\Sigma,a$}
\EndFunction
\vspace{1pt}
\Function{propose\_block}{$(H,\mathit{ac\_payload})$}
  \Comment{runs at $t_s$}
  \State $B\gets$ a block with $B.\parent=H$, $B.\slot=\slot(t)$,
    $B.\proposer=\valindex$, $B.\acp=\mathit{ac\_payload}$
  \State $B.\mathsf{attestations}\gets
    \Call{proposal\_attestations}{\Sigma,H}$
  \State broadcast $B$; \Call{on\_block}{$B$}
\EndFunction
\vspace{1pt}
\Function{attest}{$\mathit{fg\_confirmed}$}
  \State $a\gets\Call{fg\_attest}{\Sigma,\Lambda,\mathit{fg\_confirmed}}$
  \If{$a=\bot$} \State \Return \EndIf
  \State broadcast $a$; \Call{on\_attestation}{$a$}
\EndFunction
\end{algorithmic}
}

